\documentclass{amsart}
% For dealing with references we use the comment environment
\usepackage{verbatim}
\newenvironment{reference}{\comment}{\endcomment}
%\newenvironment{reference}{}{}
\newenvironment{slogan}{\comment}{\endcomment}
\newenvironment{history}{\comment}{\endcomment}

% For commutative diagrams we use Xy-pic
\usepackage[all]{xy}

% We use 2cell for 2-commutative diagrams.
\xyoption{2cell}
\UseAllTwocells

% We use multicol for the list of chapters between chapters
\usepackage{multicol}

% This is generally recommended for better output
\usepackage{lmodern}
\usepackage[T1]{fontenc}

% For cross-file-references

% Package for hypertext links:
\usepackage{hyperref}

% For any local file, say "hello.tex" you want to link to please

% Theorem environments.
%
\theoremstyle{plain}
\newtheorem{theorem}[subsection]{Theorem}
\newtheorem{proposition}[subsection]{Proposition}
\newtheorem{lemma}[subsection]{Lemma}

\theoremstyle{definition}
\newtheorem{definition}[subsection]{Definition}
\newtheorem{example}[subsection]{Example}
\newtheorem{exercise}[subsection]{Exercise}
\newtheorem{situation}[subsection]{Situation}

\theoremstyle{remark}
\newtheorem{remark}[subsection]{Remark}
\newtheorem{remarks}[subsection]{Remarks}

\numberwithin{equation}{subsection}

% Macros
%
\def\lim{\mathop{\mathrm{lim}}\nolimits}
\def\colim{\mathop{\mathrm{colim}}\nolimits}
\def\Spec{\mathop{\mathrm{Spec}}}
\def\Hom{\mathop{\mathrm{Hom}}\nolimits}
\def\Ext{\mathop{\mathrm{Ext}}\nolimits}
\def\SheafHom{\mathop{\mathcal{H}\!\mathit{om}}\nolimits}
\def\SheafExt{\mathop{\mathcal{E}\!\mathit{xt}}\nolimits}
\def\Sch{\mathit{Sch}}
\def\Mor{\mathop{\mathrm{Mor}}\nolimits}
\def\Ob{\mathop{\mathrm{Ob}}\nolimits}
\def\Sh{\mathop{\mathit{Sh}}\nolimits}
\def\NL{\mathop{N\!L}\nolimits}
\def\CH{\mathop{\mathrm{CH}}\nolimits}
\def\proetale{{pro\text{-}\acute{e}tale}}
\def\etale{{\acute{e}tale}}
\def\QCoh{\mathit{QCoh}}
\def\Ker{\mathop{\mathrm{Ker}}}
\def\Im{\mathop{\mathrm{Im}}}
\def\Coker{\mathop{\mathrm{Coker}}}
\def\Coim{\mathop{\mathrm{Coim}}}

% Boxtimes
%
\DeclareMathSymbol{\boxtimes}{\mathbin}{AMSa}{"02}

%
% Macros for moduli stacks/spaces
%
\def\QCohstack{\mathcal{QC}\!\mathit{oh}}
\def\Cohstack{\mathcal{C}\!\mathit{oh}}
\def\Spacesstack{\mathcal{S}\!\mathit{paces}}
\def\Quotfunctor{\mathrm{Quot}}
\def\Hilbfunctor{\mathrm{Hilb}}
\def\Curvesstack{\mathcal{C}\!\mathit{urves}}
\def\Polarizedstack{\mathcal{P}\!\mathit{olarized}}
\def\Complexesstack{\mathcal{C}\!\mathit{omplexes}}
% \Pic is the operator that assigns to X its picard group, usage \Pic(X)
% \Picardstack_{X/B} denotes the Picard stack of X over B
% \Picardfunctor_{X/B} denotes the Picard functor of X over B
\def\Pic{\mathop{\mathrm{Pic}}\nolimits}
\def\Picardstack{\mathcal{P}\!\mathit{ic}}
\def\Picardfunctor{\mathrm{Pic}}
\def\Deformationcategory{\mathcal{D}\!\mathit{ef}}

\setlength{\emergencystretch}{3em}
\begin{document}
\title[Stacks Project, Tag 02QD]{458 --- Periodic Euler lengths equal orders of generic periodic determinants on one-dimensional support}
\maketitle
\noindent Copyright (C) 2005--2025 Johan de Jong. Stacks Project authors. Source: \href{https://stacks.math.columbia.edu/tag/02QD}{Tag 02QD}.

\noindent Editorial difficulty note (not author text). Difficulty H3: The proof must reconcile finite-length cohomology with determinant scalars at all minimal support primes. After separating generic supports, carefully chosen good sequences in kernels and images must be lifted and adjusted to satisfy annihilator conditions, then compared through a common lattice so all length discrepancies exactly match the determinant order..

\noindent Editorial provenance note (not author text). History: excerpt from revision \texttt{a0444\allowbreak 6e57e\allowbreak c1fbc\allowbreak 252a8\allowbreak 71afc\allowbreak ec775\allowbreak 2fb28\allowbreak 07b14}, chow.tex, lines 18023--18272. Reference presentation was adapted; original-author context and core support blocks were retained or added. The mathematical wording is unchanged. Licensed under GNU FDL 1.2 or later; no invariant sections or cover texts. See ../licenses/COPYING. Context blocks reproduce author definitions and supporting results; they are not additional counted problems.

\begin{definition}
\label{definition-periodic-complex}
Let $R$ be a ring.
\begin{enumerate}
\item A {\it $2$-periodic complex} over $R$ is given
by a quadruple $(M, N, \varphi, \psi)$ consisting of
$R$-modules $M$, $N$ and $R$-module maps $\varphi : M \to N$,
$\psi : N \to M$ such that
$$
\xymatrix{
\ldots \ar[r] &
M \ar[r]^\varphi &
N \ar[r]^\psi &
M \ar[r]^\varphi &
N \ar[r] & \ldots
}
$$
is a complex. In this setting we define the {\it cohomology modules}
of the complex to be the $R$-modules
$$
H^0(M, N, \varphi, \psi) = \Ker(\varphi)/\Im(\psi)
\quad\text{and}\quad
H^1(M, N, \varphi, \psi) = \Ker(\psi)/\Im(\varphi).
$$
We say the $2$-periodic complex is {\it exact} if the cohomology
groups are zero.
\item A {\it $(2, 1)$-periodic complex} over $R$ is given
by a triple $(M, \varphi, \psi)$ consisting of an $R$-module $M$ and
$R$-module maps $\varphi : M \to M$, $\psi : M \to M$
such that
$$
\xymatrix{
\ldots \ar[r] &
M \ar[r]^\varphi &
M \ar[r]^\psi &
M \ar[r]^\varphi &
M \ar[r] & \ldots
}
$$
is a complex. Since this is a special case of a $2$-periodic complex
we have its {\it cohomology modules} $H^0(M, \varphi, \psi)$,
$H^1(M, \varphi, \psi)$ and a notion of exactness.
\end{enumerate}
\end{definition}

\begin{definition}
\label{definition-periodic-length}
Let $(M, N, \varphi, \psi)$ be a $2$-periodic complex
over a ring $R$ whose cohomology modules have finite length.
In this case we define the {\it multiplicity} of $(M, N, \varphi, \psi)$
to be the integer
$$
e_R(M, N, \varphi, \psi) =
\text{length}_R(H^0(M, N, \varphi, \psi))
-
\text{length}_R(H^1(M, N, \varphi, \psi))
$$
In the case of a $(2, 1)$-periodic complex $(M, \varphi, \psi)$,
we denote this by $e_R(M, \varphi, \psi)$ and we will sometimes call this
the {\it (additive) Herbrand quotient}.
\end{definition}

\section{Appendix A: Alternative approach to key lemma}
\label{section-appendix-A}

\noindent
In this appendix we first define determinants $\det_\kappa(M)$
of finite length modules $M$ over local rings $(R, \mathfrak m, \kappa)$,
see Subsection \href{https://stacks.math.columbia.edu/tag/02P5}{subsection determinants finite length (Tag 02P5)}.
The determinant $\det_\kappa(M)$ is a $1$-dimensional $\kappa$-vector space.
We use this in Subsection \href{https://stacks.math.columbia.edu/tag/0BDR}{subsection periodic complexes determinants (Tag 0BDR)}
to define the determinant $\det_\kappa(M, \varphi, \psi) \in \kappa^*$
of an exact $(2, 1)$-periodic complex $(M, \varphi, \psi)$
with $M$ of finite length. In Subsection \href{https://stacks.math.columbia.edu/tag/02PW}{subsection symbols (Tag 02PW)}
we use these determinants to construct a tame symbol
$d_R(a, b) = \det_\kappa(R/ab, a, b)$ for a pair of nonzerodivisors
$a, b \in R$ when $R$ is Noetherian of dimension $1$.
Although there is no doubt that
$$
d_R(a, b) = \partial_R(a, b)
$$
where $\partial_R$ is as in Section \href{https://stacks.math.columbia.edu/tag/0EAH}{section tame symbol (Tag 0EAH)},
we have not (yet) added the verification. The advantage of the
tame symbol as constructed in this appendix is that it extends
(for example) to pairs of injective endomorphisms $\varphi, \psi$
of a finite $R$-module $M$ of dimension $1$ such that
$\varphi(\psi(M)) = \psi(\varphi(M))$. In
Subsection \href{https://stacks.math.columbia.edu/tag/02QA}{subsection length determinant (Tag 02QA)}
we relate Herbrand quotients and determinants.
An easy to state version of the main result
(Proposition \ref{proposition-length-determinant-periodic-complex})
is the formula
$$
-e_R(M, \varphi, \psi) =
\text{ord}_R(\det\nolimits_K(M_K, \varphi, \psi))
$$
when $(M, \varphi, \psi)$ is a $(2, 1)$-periodic complex
whose Herbrand quotient $e_R$ (Definition \ref{definition-periodic-length})
is defined
over a $1$-dimensonal Noetherian local domain $R$ with fraction field $K$.
We use this proposition to give an alternative proof of the key lemma
(Lemma \href{https://stacks.math.columbia.edu/tag/0EAX}{lemma milnor gersten low degree (Tag 0EAX)})
for the tame symbol constructed in this appendix, see
Lemma \href{https://stacks.math.columbia.edu/tag/02QJ}{lemma secondary ramification (Tag 02QJ)}.




\begin{definition}
\label{definition-periodic-determinant}
Let $R$ be a local ring with residue field $\kappa$.
Let $(M, \varphi, \psi)$ be a $(2, 1)$-periodic complex over $R$.
Assume that $M$ has finite length and that $(M, \varphi, \psi)$ is
exact. The {\it determinant of $(M, \varphi, \psi)$} is
the element
$$
\det\nolimits_\kappa(M, \varphi, \psi) \in \kappa^*
$$
such that the composition
$$
\det\nolimits_\kappa(M)
\xrightarrow{\gamma_\psi \circ \sigma \circ \gamma_\varphi^{-1}}
\det\nolimits_\kappa(M)
$$
is multiplication by
$(-1)^{\text{length}_R(I_\varphi)\text{length}_R(I_\psi)}
\det\nolimits_\kappa(M, \varphi, \psi)$.
\end{definition}

\begin{definition}
\label{definition-determinant}
Let $R$ be a local ring with maximal ideal $\mathfrak m$ and
residue field $\kappa$. Let $M$ be a finite length $R$-module.
Say $l = \text{length}_R(M)$.
\begin{enumerate}
\item Given elements $x_1, \ldots, x_r \in M$ we denote
$\langle x_1, \ldots, x_r \rangle = Rx_1 + \ldots + Rx_r$ the
$R$-submodule of $M$ generated by $x_1, \ldots, x_r$.
\item We will say an $l$-tuple of elements
$(e_1, \ldots, e_l)$ of $M$ is {\it admissible} if
$\mathfrak m e_i \subset \langle e_1, \ldots, e_{i - 1} \rangle$
for $i = 1, \ldots, l$.
\item A {\it symbol} $[e_1, \ldots, e_l]$ will mean
$(e_1, \ldots, e_l)$ is an admissible $l$-tuple.
\item An {\it admissible relation} between symbols is one of the following:
\begin{enumerate}
\item if $(e_1, \ldots, e_l)$ is an admissible sequence and
for some $1 \leq a \leq l$ we have
$e_a \in \langle e_1, \ldots, e_{a - 1}\rangle$, then
$[e_1, \ldots, e_l] = 0$,
\item if $(e_1, \ldots, e_l)$ is an admissible sequence and
for some $1 \leq a \leq l$ we have $e_a = \lambda e'_a + x$
with $\lambda \in R^*$, and
$x \in \langle e_1, \ldots, e_{a - 1}\rangle$, then
$$
[e_1, \ldots, e_l] =
\overline{\lambda} [e_1, \ldots, e_{a - 1}, e'_a, e_{a + 1}, \ldots, e_l]
$$
where $\overline{\lambda} \in \kappa^*$ is the image of $\lambda$ in
the residue field, and
\item if $(e_1, \ldots, e_l)$ is an admissible sequence and
$\mathfrak m e_a \subset \langle e_1, \ldots, e_{a - 2}\rangle$ then
$$
[e_1, \ldots, e_l] =
- [e_1, \ldots, e_{a - 2}, e_a, e_{a - 1}, e_{a + 1}, \ldots, e_l].
$$
\end{enumerate}
\item
We define the {\it determinant of the finite length $R$-module $M$} to be
$$
\det\nolimits_\kappa(M) =
\left\{
\frac{\kappa\text{-vector space generated by symbols}}
{\kappa\text{-linear combinations of admissible relations}}
\right\}
$$
\end{enumerate}
\end{definition}

\noindent
Let $R$ be a local ring with residue field $\kappa$.
Let $(M, \varphi, \psi)$ be a $(2, 1)$-periodic complex over $R$.
Assume that $M$ has finite length and that $(M, \varphi, \psi)$ is
exact. We are going to use the determinant construction to define
an invariant of this situation. See
Subsection \href{https://stacks.math.columbia.edu/tag/02P5}{subsection determinants finite length (Tag 02P5)}.
Let us abbreviate
$K_\varphi = \Ker(\varphi)$,
$I_\varphi = \Im(\varphi)$,
$K_\psi = \Ker(\psi)$, and
$I_\psi = \Im(\psi)$.
The short exact sequences
$$
0 \to K_\varphi \to M \to I_\varphi \to 0, \quad
0 \to K_\psi \to M \to I_\psi \to 0
$$
give isomorphisms
$$
\gamma_\varphi :
\det\nolimits_\kappa(K_\varphi)
\otimes
\det\nolimits_\kappa(I_\varphi)
\longrightarrow
\det\nolimits_\kappa(M), \quad
\gamma_\psi :
\det\nolimits_\kappa(K_\psi)
\otimes
\det\nolimits_\kappa(I_\psi)
\longrightarrow
\det\nolimits_\kappa(M),
$$
see Lemma \href{https://stacks.math.columbia.edu/tag/02PA}{lemma det exact sequences (Tag 02PA)}.
On the other hand the exactness of the complex gives equalities
$K_\varphi = I_\psi$, and $K_\psi = I_\varphi$
and hence an isomorphism
$$
\sigma :
\det\nolimits_\kappa(K_\varphi)
\otimes
\det\nolimits_\kappa(I_\varphi)
\longrightarrow
\det\nolimits_\kappa(K_\psi)
\otimes
\det\nolimits_\kappa(I_\psi)
$$
by switching the factors. Using this notation we can define our invariant.


\begin{proposition}
\label{proposition-length-determinant-periodic-complex}
Let $R$ be a local Noetherian ring with residue field $\kappa$.
Suppose that $(M, \varphi, \psi)$ is a $(2, 1)$-periodic
complex over $R$. Assume
\begin{enumerate}
\item $M$ is a finite $R$-module,
\item the cohomology modules of $(M, \varphi, \psi)$ are of finite length, and
\item $\dim(\text{Supp}(M)) = 1$.
\end{enumerate}
Let $\mathfrak q_i$, $i = 1, \ldots, t$ be the minimal
primes of the support of $M$. Then we have\footnote{
Obviously we could get rid of the minus sign by redefining
$\det_\kappa(M, \varphi, \psi)$ as the inverse of its
current value, see Definition \ref{definition-periodic-determinant}.}
$$
- e_R(M, \varphi, \psi) =
\sum\nolimits_{i = 1, \ldots, t}
\text{ord}_{R/\mathfrak q_i}\left(
\det\nolimits_{\kappa(\mathfrak q_i)}
(M_{\mathfrak q_i}, \varphi_{\mathfrak q_i}, \psi_{\mathfrak q_i})
\right)
$$
\end{proposition}

\begin{proof}
We first reduce to the case $t = 1$ in the following way.
Note that
$\text{Supp}(M) = \{\mathfrak m, \mathfrak q_1, \ldots, \mathfrak q_t\}$,
where $\mathfrak m \subset R$ is the maximal ideal.
Let $M_i$ denote the image of $M \to M_{\mathfrak q_i}$,
so $\text{Supp}(M_i) = \{\mathfrak m, \mathfrak q_i\}$.
The map $\varphi$ (resp.\ $\psi$) induces an $R$-module map
$\varphi_i : M_i \to M_i$ (resp.\ $\psi_i : M_i \to M_i$).
Thus we get a morphism of $(2, 1)$-periodic complexes
$$
(M, \varphi, \psi) \longrightarrow
\bigoplus\nolimits_{i = 1, \ldots, t} (M_i, \varphi_i, \psi_i).
$$
The kernel and cokernel of this map have support contained in
$\{\mathfrak m\}$. Hence by Lemma \href{https://stacks.math.columbia.edu/tag/0EA9}{lemma compare periodic lengths (Tag 0EA9)}
we have
$$
e_R(M, \varphi, \psi) =
\sum\nolimits_{i = 1, \ldots, t}
e_R(M_i, \varphi_i, \psi_i)
$$
On the other hand we clearly have $M_{\mathfrak q_i} = M_{i, \mathfrak q_i}$,
and hence the terms of the right hand side of the formula of the
lemma are equal to the expressions
$$
\text{ord}_{R/\mathfrak q_i}\left(
\det\nolimits_{\kappa(\mathfrak q_i)}
(M_{i, \mathfrak q_i}, \varphi_{i, \mathfrak q_i}, \psi_{i, \mathfrak q_i})
\right)
$$
In other words, if we can prove the lemma for each of the modules
$M_i$, then the lemma holds. This reduces us to the case $t = 1$.

\medskip\noindent
Assume we have a $(2, 1)$-periodic complex $(M, \varphi, \psi)$
over a Noetherian local ring with $M$ a finite $R$-module,
$\text{Supp}(M) = \{\mathfrak m, \mathfrak q\}$, and
finite length cohomology modules. The proof in this case
follows from Lemma \href{https://stacks.math.columbia.edu/tag/02QB}{lemma key lemma (Tag 02QB)} and careful bookkeeping.
Denote
$K_\varphi = \Ker(\varphi)$,
$I_\varphi = \Im(\varphi)$,
$K_\psi = \Ker(\psi)$, and
$I_\psi = \Im(\psi)$.
Since $R$ is Noetherian these are all finite $R$-modules.
Set
$$
a = \text{length}_{R_{\mathfrak q}}(I_{\varphi, \mathfrak q})
= \text{length}_{R_{\mathfrak q}}(K_{\psi, \mathfrak q}),
\quad
b = \text{length}_{R_{\mathfrak q}}(I_{\psi, \mathfrak q})
= \text{length}_{R_{\mathfrak q}}(K_{\varphi, \mathfrak q}).
$$
Equalities because the complex becomes exact after localizing at
$\mathfrak q$. Note that $l = \text{length}_{R_{\mathfrak q}}(M_{\mathfrak q})$
is equal to $l = a + b$.

\medskip\noindent
We are going to use Lemma \href{https://stacks.math.columbia.edu/tag/02QC}{lemma good sequence exists (Tag 02QC)}
to choose sequences of elements in finite $R$-modules
$N$ with support contained in $\{\mathfrak m, \mathfrak q\}$.
In this case $N_{\mathfrak q}$ has finite length, say $n \in \mathbf{N}$.
Let us call a sequence $w_1, \ldots, w_n \in N$
with properties (1) and (2) of Lemma \href{https://stacks.math.columbia.edu/tag/02QC}{lemma good sequence exists (Tag 02QC)}
a ``good sequence''. Note that the quotient
$N/\langle w_1, \ldots, w_n \rangle$ of $N$ by the submodule generated by
a good sequence has support (contained in) $\{\mathfrak m\}$
and hence has finite length (Algebra, Lemma \href{https://stacks.math.columbia.edu/tag/00L5}{lemma support point (Tag 00L5)}).
Moreover, the symbol
$[w_1, \ldots, w_n] \in \det_{\kappa(\mathfrak q)}(N_{\mathfrak q})$
is a generator, see Lemma \href{https://stacks.math.columbia.edu/tag/02P9}{lemma determinant dimension one (Tag 02P9)}.

\medskip\noindent
Having said this we choose good sequences
$$
\begin{matrix}
x_1, \ldots, x_b & \text{in} & K_\varphi, &
t_1, \ldots, t_a & \text{in} & K_\psi, \\
y_1, \ldots, y_a & \text{in} & I_\varphi \cap \langle t_1, \ldots t_a\rangle, &
s_1, \ldots, s_b & \text{in} & I_\psi \cap \langle x_1, \ldots, x_b\rangle.
\end{matrix}
$$
We will adjust our choices a little bit as follows.
Choose lifts $\tilde y_i \in M$ of $y_i \in I_\varphi$
and $\tilde s_i \in M$ of $s_i \in I_\psi$. It may not be the case
that $\mathfrak q \tilde y_1 \subset \langle x_1, \ldots, x_b\rangle$
and it may not be the case that
$\mathfrak q \tilde s_1 \subset \langle t_1, \ldots, t_a\rangle$.
However, using that $\mathfrak q$ is finitely generated (as in the proof
of Lemma \href{https://stacks.math.columbia.edu/tag/02QC}{lemma good sequence exists (Tag 02QC)}) we can find a
$d \in R$, $d \not \in \mathfrak q$ such that
$\mathfrak q d\tilde y_1 \subset \langle x_1, \ldots, x_b\rangle$
and
$\mathfrak q d\tilde s_1 \subset \langle t_1, \ldots, t_a\rangle$.
Thus after replacing $y_i$ by $dy_i$,
$\tilde y_i$ by $d\tilde y_i$, $s_i$ by $ds_i$ and $\tilde s_i$
by $d\tilde s_i$ we see that we may assume also that
$x_1, \ldots, x_b, \tilde y_1, \ldots, \tilde y_b$
and $t_1, \ldots, t_a, \tilde s_1, \ldots, \tilde s_b$
are good sequences in $M$.

\medskip\noindent
Finally, we choose a good sequence
$z_1, \ldots, z_l$ in the finite $R$-module
$$
\langle
x_1, \ldots, x_b, \tilde y_1, \ldots, \tilde y_a
\rangle
\cap
\langle
t_1, \ldots, t_a, \tilde s_1, \ldots, \tilde s_b
\rangle.
$$
Note that this is also a good sequence in $M$.

\medskip\noindent
Since $I_{\varphi, \mathfrak q} = K_{\psi, \mathfrak q}$
there is a unique element $h \in \kappa(\mathfrak q)$ such that
$[y_1, \ldots, y_a] = h [t_1, \ldots, t_a]$
inside $\det_{\kappa(\mathfrak q)}(K_{\psi, \mathfrak q})$.
Similarly, as $I_{\psi, \mathfrak q} = K_{\varphi, \mathfrak q}$
there is a unique element $h \in \kappa(\mathfrak q)$ such that
$[s_1, \ldots, s_b] = g [x_1, \ldots, x_b]$
inside $\det_{\kappa(\mathfrak q)}(K_{\varphi, \mathfrak q})$.
We can also do this with the three good sequences we have
in $M$. All in all we get the following identities
\begin{align*}
[y_1, \ldots, y_a]
& =
h [t_1, \ldots, t_a] \\
[s_1, \ldots, s_b]
& =
g [x_1, \ldots, x_b] \\
[z_1, \ldots, z_l]
& =
f_\varphi [x_1, \ldots, x_b, \tilde y_1, \ldots, \tilde y_a] \\
[z_1, \ldots, z_l]
& =
f_\psi [t_1, \ldots, t_a, \tilde s_1, \ldots, \tilde s_b]
\end{align*}
for some $g, h, f_\varphi, f_\psi \in \kappa(\mathfrak q)$.

\medskip\noindent
Having set up all this
notation let us compute $\det_{\kappa(\mathfrak q)}(M, \varphi, \psi)$.
Namely, consider the element $[z_1, \ldots, z_l]$.
Under the map $\gamma_\psi \circ \sigma \circ \gamma_\varphi^{-1}$
of Definition \ref{definition-periodic-determinant} we have
\begin{eqnarray*}
[z_1, \ldots, z_l] & = &
f_\varphi [x_1, \ldots, x_b, \tilde y_1, \ldots, \tilde y_a] \\
& \mapsto & f_\varphi [x_1, \ldots, x_b] \otimes [y_1, \ldots, y_a] \\
& \mapsto &
f_\varphi h/g [t_1, \ldots, t_a] \otimes [s_1, \ldots, s_b] \\
& \mapsto &
f_\varphi h/g [t_1, \ldots, t_a, \tilde s_1, \ldots, \tilde s_b] \\
& = &
f_\varphi h/f_\psi g [z_1, \ldots, z_l]
\end{eqnarray*}
This means that
$\det_{\kappa(\mathfrak q)}
(M_{\mathfrak q}, \varphi_{\mathfrak q}, \psi_{\mathfrak q})$
is equal to $f_\varphi h/f_\psi g$ up to a sign.

\medskip\noindent
We abbreviate the following quantities
\begin{eqnarray*}
k_\varphi & = & \text{length}_R(K_\varphi/\langle x_1, \ldots, x_b\rangle) \\
k_\psi    & = & \text{length}_R(K_\psi/\langle t_1, \ldots, t_a\rangle) \\
i_\varphi & = & \text{length}_R(I_\varphi/\langle y_1, \ldots, y_a\rangle) \\
i_\psi    & = & \text{length}_R(I_\psi/\langle s_1, \ldots, s_a\rangle) \\
m_\varphi & = & \text{length}_R(M/
\langle x_1, \ldots, x_b, \tilde y_1, \ldots, \tilde y_a\rangle) \\
m_\psi    & = & \text{length}_R(M/
\langle t_1, \ldots, t_a, \tilde s_1, \ldots, \tilde s_b\rangle) \\
\delta_\varphi & = & \text{length}_R(
\langle x_1, \ldots, x_b, \tilde y_1, \ldots, \tilde y_a\rangle
\langle z_1, \ldots, z_l\rangle) \\
\delta_\psi & = & \text{length}_R(
\langle t_1, \ldots, t_a, \tilde s_1, \ldots, \tilde s_b\rangle
\langle z_1, \ldots, z_l\rangle)
\end{eqnarray*}
Using the exact sequences $0 \to K_\varphi \to M \to I_\varphi \to 0$
we get $m_\varphi = k_\varphi + i_\varphi$. Similarly we have
$m_\psi = k_\psi + i_\psi$. We have
$\delta_\varphi + m_\varphi = \delta_\psi + m_\psi$ since this
is equal to the colength of $\langle z_1, \ldots, z_l \rangle$
in $M$. Finally, we have
$$
\delta_\varphi = \text{ord}_{R/\mathfrak q}(f_\varphi),
\quad
\delta_\psi = \text{ord}_{R/\mathfrak q}(f_\psi)
$$
by our first application of the key Lemma \href{https://stacks.math.columbia.edu/tag/02QB}{lemma key lemma (Tag 02QB)}.

\medskip\noindent
Next, let us compute the multiplicity of the periodic complex
\begin{eqnarray*}
e_R(M, \varphi, \psi) & = &
\text{length}_R(K_\varphi/I_\psi) - \text{length}_R(K_\psi/I_\varphi) \\
& = &
\text{length}_R(
\langle x_1, \ldots, x_b\rangle/
\langle s_1, \ldots, s_b\rangle)
+ k_\varphi - i_\psi \\
& & -
\text{length}_R(
\langle t_1, \ldots, t_a\rangle/
\langle y_1, \ldots, y_a\rangle)
- k_\psi + i_\varphi \\
& = &
\text{ord}_{R/\mathfrak q}(g/h) + k_\varphi - i_\psi - k_\psi + i_\varphi \\
& = &
\text{ord}_{R/\mathfrak q}(g/h) + m_\varphi - m_\psi \\
& = &
\text{ord}_{R/\mathfrak q}(g/h) + \delta_\psi - \delta_\varphi \\
& = &
\text{ord}_{R/\mathfrak q}(f_\psi g/f_\varphi h)
\end{eqnarray*}
where we used the key Lemma \href{https://stacks.math.columbia.edu/tag/02QB}{lemma key lemma (Tag 02QB)} twice in the third equality.
By our computation of $\det_{\kappa(\mathfrak q)}
(M_{\mathfrak q}, \varphi_{\mathfrak q}, \psi_{\mathfrak q})$
this proves the proposition.
\end{proof}
\end{document}
